\documentclass[11pt]{article}
\usepackage[a4paper,margin=20mm]{geometry}
\usepackage{fontspec}
\setmainfont{Noto Serif CJK SC}
\XeTeXlinebreaklocale "zh"
\XeTeXlinebreakskip=0pt plus 1pt
\usepackage{amsmath,amssymb,amsthm,mathtools,tikz-cd,enumerate,paralist,xurl,hyperref}
\setlength{\emergencystretch}{4em}
\newcommand{\identity}{\mathrm{id}}
\newcommand{\Obj}{\operatorname{Ob}}
\newcommand{\Mor}{\operatorname{Mor}}
\newcommand{\Hom}{\operatorname{Hom}}
\newcommand{\End}{\operatorname{End}}
\newcommand{\Ker}{\operatorname{ker}}
\newcommand{\Coker}{\operatorname{coker}}
\newcommand{\opp}{\mathrm{op}}
\newcommand{\munit}{\mathbf{1}}
\newcommand{\Z}{\mathbb{Z}}
\newcommand{\cate}[1]{\mathbf{#1}}
\newcommand{\cated}[1]{\mathbf{#1}\text{-}}
\newcommand{\dcate}[1]{\text{-}\mathbf{#1}}
\newcommand{\dotimes}[1]{\mathbin{\mathop{\otimes}\limits_{#1}}}
\newcommand{\rightiso}{\xrightarrow{\sim}}
\renewcommand{\index}[2][]{\relax}
\newcommand{\CorpusNumber}{0}
\newtheorem{theorem}{定理}
\newtheorem{lemma}[theorem]{引理}
\newtheorem{proposition}[theorem]{命题}
\newtheorem{corollary}[theorem]{推论}
\theoremstyle{definition}
\newtheorem{definition}[theorem]{定义}
\newtheorem{definition-proposition}[theorem]{定义--命题}
\newtheorem{example}[theorem]{例}
\newtheorem{remark}[theorem]{注记}
\renewcommand{\thetheorem}{\CorpusNumber}
\renewcommand{\proofname}{证明}
\begin{document}
\section*{295: Gabriel compact-projective-generator theorem}
\noindent\textbf{作者：} 李文威 (Wen-Wei Li). \textbf{作品：}《代数学方法》第二卷，第七章。
\par\noindent\textbf{来源与版本：} \url{https://github.com/wenweili/AlJabr-2/blob/e03224e2e867a49dc60c7975d890079068a30466/chapter7.tex#L2071}, commit \texttt{e03224e2e867a49dc60c7975d890079068a30466}.
\par\noindent\textbf{许可：} Creative Commons Attribution 4.0 International; copyright 2024 李文威. \url{https://creativecommons.org/licenses/by/4.0/}. 作者不保证无误；见附带 LICENSE。
\par\noindent\textbf{整理说明（非作者正文）：} 下列题面、证明及局部支撑正文来自作者原生 TeX；保留原语言、顺序及数学内容。仅改独立排版、原编号引用与索引命令；未生成或补写证明。所有文字与数学图表均可编辑。标准先修仅引用，不递归重证。
\par\noindent\textbf{难度说明（整理者）：} H2: The proof reconstructs arbitrary modules from presentations using a compact projective generator, and proves fullness and essential surjectivity of the Hom functor. This category reconstruction exceeds routine calculations with projective modules.
\subsection*{原命题与完整人类证明}
\def\CorpusNumber{7.8.4}
\begin{theorem}[P.\ Gabriel]\label{prop:identify-ModR}
	设 $s \in \Obj(\mathcal{A})$, 命 $R := \End_{\mathcal{A}}(s)$, 则
	\[ \Hom_{\mathcal{A}}(s, \cdot): \mathcal{A} \to \cated{Mod}R \]
	为等价的充要条件是 $s$ 为 $\mathcal{A}$ 的紧投射生成元.
\end{theorem}
\begin{proof}
	简记 $G := \Hom_{\mathcal{A}}(s, \cdot): \mathcal{A} \to \cated{Mod}R$. 先说明条件的必要性. 由于 $G$ 是等价并且 $G(s) = R$, 问题归结为验证 $R$ 是 $\cated{Mod}R$ 的紧投射生成元, 这是简单的.
	
	以下处理充分性. 设 $s$ 是紧投射生成元. 命题 7.8.3 蕴涵 $G$ 是保所有小 $\varinjlim$ 的忠实正合函子, 因为它与忘却 $\cated{Mod}R \to \Bbbk\dcate{Mod}$ 的合成有此性质. 既然忠实性已知, 问题归结为证:
	\begin{compactitem}
		\item 对所有 $X, Y \in \Obj(\mathcal{A})$, 相应的 $\Hom_{\mathcal{A}}(X, Y) \to \Hom_R(GX, GY)$ 是满射;
		\item $G$ 是本质满的.
	\end{compactitem}
	
	我们首先说明对于任何 $X \in \Obj(\mathcal{A})$, 定义态射 $\phi_X: s^{\oplus \Hom_{\mathcal{A}}(s, X)} \to X$ 使得它在 $f: s \to X$ 对应的直和项上是 $f$, 则 $\phi_X$ 是满的. 使用反证法: 若 $\Coker(\phi_X) \neq 0$, 则生成元的性质说明存在非零态射 $s \to \Coker(\phi_X)$ (命题 2.8.15), 而投射对象的性质说明此态射可以分解为 $s \xrightarrow{g} X \to \Coker(\phi_X)$, 故 $s^{\oplus \Hom_{\mathcal{A}}(s, X)} \xrightarrow{\phi_X} X \to \Coker(\phi_X)$ 在 $g$ 对应的直和项上非零, 矛盾.
	
	对 $\Ker(\phi_X)$ 的核也可以继续如上操作, 由此得知任何 $X$ 皆可置入正合列
	\[ s^{\oplus J} \to s^{\oplus I} \to X \to 0, \quad I = I_X, \; J = J_X: \text{小集}. \]
	因为 $G$ 保所有小 $\varinjlim$, 这给出右 $R$-模的正合列
	\[ R^{\oplus J} \to R^{\oplus I} \to GX \to 0, \]
	两者的第一段态射都可以视同 $R$ 上的 $I \times J$ 矩阵, 以左乘作用; 它们实际是同一个矩阵.
	
	现在考虑右 $R$-模同态 $f: GX \to GY$. 对 $X$ 和 $Y$ 任取如上形式的正合列, 易见 $f$ 可提升为行正合交换图表
	\[\begin{tikzcd}
		R^{\oplus J_X} \arrow[r] \arrow[d] & R^{\oplus I_X} \arrow[d] \arrow[r] & GX \arrow[d, "f"] \arrow[r] & 0 \\
		R^{\oplus J_Y} \arrow[r] & R^{\oplus I_Y} \arrow[r] & GY \arrow[r] & 0.
	\end{tikzcd}\]
	但左侧方块的四个箭头都可以表成 $R$ 上的矩阵, 四个矩阵遂给出行正合交换图表
	\[\begin{tikzcd}
		s^{\oplus J_X} \arrow[r] \arrow[d] & s^{\oplus I_X} \arrow[d] \arrow[r] & X \arrow[dashed, d] \arrow[r] & 0 \\
		s^{\oplus J_Y} \arrow[r] & s^{\oplus I_Y} \arrow[r] & Y \arrow[r] & 0
	\end{tikzcd}\]
	的实线部分, 虚线部分由余核的函子性唯一确定. 由于 $G$ 正合, 必有 $G(X \dashrightarrow Y) = f$.
	
	类似的技巧可以说明 $G$ 本质满. 给定右 $R$-模 $M$, 存在正合列
	\[ R^{\oplus L} \to R^{\oplus K} \to M \to 0, \]
	其第一段映射视同 $R$ 上的 $K \times L$ 矩阵. 但同一个矩阵在 $\mathcal{A}$ 中给出余核正合列
	\[ s^{\oplus L} \to s^{\oplus K} \to X \to 0; \]
	既然 $G$ 保所有小 $\varinjlim$, 必有 $GX \simeq M$. 明所欲证.
\end{proof}
\subsection*{必要作者背景与局部支撑（不另计题）}
\par\medskip\noindent\textit{作者原文：chapter7.tex，行 2037--2069.}\par
本节依然选定交换环 $\Bbbk$. 所有 Abel 范畴和函子皆默认是 $\Bbbk$-线性的.

形如 $R\dcate{Mod}$ 或 $\cated{Mod}R$ 的范畴统称为模范畴, 其中 $R$ 是 $\Bbbk$-代数. 我们在 \S7.7 对模范畴与其间的函子进行了专门的讨论. 另一方面, 当然也可以问有哪些范畴等价于模范畴, 以及等价的具体给法. 本节将顺势给出几则简短而方便的相关结果.

设 $\mathcal{A}$ 为 Abel 范畴. 本节的进路是考虑形如 $\Hom_{\mathcal{A}}(s, \cdot)$ 的函子, 它自然地升级为函子 $\mathcal{A} \to \cated{Mod}R$, 其中 $R := \End_{\mathcal{A}}(s)$, 问题在于判断 $\Hom_{\mathcal{A}}(s, \cdot)$ 何时为等价.

\index{jinduixiang}
兹复述紧对象的定义 A.2.2: 设 $\mathcal{A}$ 有所有滤过小 $\varinjlim$ 而 $X \in \Obj(\mathcal{A})$, 假若对所有滤过小范畴 $I$ 和函子 $\alpha: I \to \mathcal{A}$, 以下典范态射皆为同构, 则称 $X$ 紧:
\[ \varinjlim_i \Hom\left( X, \alpha(i) \right) \to \Hom\left( X, \varinjlim \alpha \right). \]

\def\CorpusNumber{7.8.1}
\begin{lemma}\label{prop:proj-compact}
	设 $X$ 是余完备 Abel 范畴 $\mathcal{A}$ 的投射对象, 则 $X$ 紧当且仅当 $\Hom_{\mathcal{A}}(X, \cdot)$ 保所有小直和, 亦即: 对任意小集 $I$ 和一族对象 $(Y_i)_{i \in I}$, 我们有 $\bigoplus_{i \in I} \Hom_{\mathcal{A}}(X, Y_i) \rightiso \Hom_{\mathcal{A}}(X, \bigoplus_{i \in I} Y_i)$. 当前述条件成立时, $\Hom_{\mathcal{A}}(X, \cdot)$ 实际上保所有小 $\varinjlim$.
\end{lemma}
\begin{proof}
	设 $X$ 紧. 因为以小集 $I$ 为下标的直和是有限直和的滤过 $\varinjlim$ (取遍 $I$ 的有限子集, 以包含关系赋序), 已知 $\Hom_{\mathcal{A}}(X, \cdot)$ 保有限直和, 故它也保小直和.
	
	对于反方向, 投射条件相当于说 $\Hom_{\mathcal{A}}(X, \cdot)$ 保余核, 而一般的小 $\varinjlim$ 可以用小直和连同余核来构造, 故 $\Hom_{\mathcal{A}}(X, \cdot)$ 保所有小 $\varinjlim$.
\end{proof}

\def\CorpusNumber{7.8.2}
\begin{example}
	设 $R$ 为 $\Bbbk$-代数, $M$ 为投射右 $R$-模, 则 $M$ 紧的充要条件是 $M$ 有限生成.
	\begin{itemize}
		\item 必要性: 存在小集 $I$, 同态 $i: M \to R^{\oplus I}$ 和 $p: R^{\oplus I} \to M$ 使得 $pi = \identity_M$. 紧性蕴涵 $i$ 取值在 $R^{\oplus I_0}$ 中, 其中 $I_0 \subset I$ 是有限子集; 取 $p' := p|_{R^{\oplus I_0}}$, 则仍有 $p' i = \identity_M$, 这就说明 $M$ 可以实现为有限秩自由模的直和项.
		\item 充分性: 代入引理 7.8.1 的判准. 任意同态 $M \to \bigoplus_{i \in I} Y_i$ 都必然落在 $\bigoplus_{i \in I_0} Y_i$ 中, 其中 $I_0 \subset I$ 是有限子集 (考察生成元的像即可), 故 $M$ 紧.
	\end{itemize}
\end{example}

\def\CorpusNumber{7.8.3}
\begin{proposition}\label{prop:proj-compact-gen}
	设 $\mathcal{A}$ 是余完备 Abel 范畴, 则 $s \in \Obj(\mathcal{A})$ 是紧投射生成元的充要条件是 $\Hom_{\mathcal{A}}(s, \cdot): \mathcal{A} \to \Bbbk\dcate{Mod}$ 是保所有小 $\varinjlim$ 的忠实正合函子.
\end{proposition}
\begin{proof}
	根据命题 2.8.15, 投射生成元等价于 $\Hom_{\mathcal{A}}(s, \cdot)$ 忠实正合. 根据引理 7.8.1, 在投射的前提下, 紧性等价于 $\Hom_{\mathcal{A}}(s, \cdot)$ 保所有小 $\varinjlim$.
\end{proof}
\subsection*{标准先修与保留说明（整理者）}
通常的伴随、锥与（余）极限、Abel 范畴、投射对象和生成元定义为标准先修。原书编号引用在本文件中保持原编号。森田分类中 Eilenberg--Watts 定理 7.7.2 为具名标准先修；本文件保留其陈述，以及实际用到的有限投射对偶、评价与余评价构造。原书命题 2.8.15 为投射生成元的忠实正合刻画；A.2.2 为紧对象定义（已在局部背景中复述）。原文的简短例行验证及任何已存在的笔误均原样保留，未作数学修改。
\begin{thebibliography}{Li1}
\bibitem{Li1} 李文威，《代数学方法》第一卷：基础架构，高等教育出版社。原书引用保留原章节和结果编号；\url{https://www.wwli.asia/mathematics/AlJabr/}.
\end{thebibliography}
\end{document}
